% META: {"id": "1SPE-EXPONENTIELLE-QCM", "chapitre": "1SPE-EXPONENTIELLE", "type_objet": "qcm", "genere_depuis": "chapitres/1SPE-EXPONENTIELLE/qcm/1SPE-EXPONENTIELLE-QCM.json", "status": "generated"}
% Fichier genere par scripts/build_qcm_tex.py — ne pas editer a la main.

\section*{\textcolor{chapcolor}{\MakeUppercase{Fonction exponentielle — QCM d'auto-évaluation}}}

\begin{center}
\textit{Pour chaque question, une seule réponse est exacte.}
\end{center}

\begin{enumerate}

\item[] \textbf{Capacité C1}

\item \textbf{[Q1]} Parmi les propriétés suivantes, laquelle caractérise la fonction exponentielle ?
  \begin{enumerate}[label=\Alph*.]
    \item C'est la seule fonction dérivable sur $\mathbb{R}$ telle que $f'(x) = f(x)$ pour tout $x$ et $f(0) = 1$.
    \item C'est la seule fonction dérivable sur $\mathbb{R}$ telle que $f'(x) = f(x)$ pour tout $x$ et $f(1) = 1$.
    \item C'est la seule fonction dérivable sur $\mathbb{R}$ telle que $f'(x) = f(x)$ pour tout $x$ et $f(0) = 0$.
    \item C'est la seule fonction dérivable sur $\mathbb{R}$ telle que $f'(x) = x f(x)$ pour tout $x$ et $f(0) = 1$.
  \end{enumerate}

\item \textbf{[Q2]} Quelle est la valeur de $\mathrm{e}^0$ ?
  \begin{enumerate}[label=\Alph*.]
    \item $0$
    \item $1$
    \item $\mathrm{e}$
    \item La valeur n'est pas définie.
  \end{enumerate}

\item \textbf{[Q3]} Pour tout réel $x$, on a :
  \begin{enumerate}[label=\Alph*.]
    \item $\mathrm{e}^x$ peut être négatif si $x < 0$
    \item $\mathrm{e}^x \geqslant 0$ (avec $\mathrm{e}^x = 0$ possible)
    \item $\mathrm{e}^x > 0$
    \item $\mathrm{e}^x > 1$ pour tout $x \in \mathbb{R}$
  \end{enumerate}

\item[] \textbf{Capacité C2}

\item \textbf{[Q4]} Simplifier $\mathrm{e}^3 \times \mathrm{e}^5$.
  \begin{enumerate}[label=\Alph*.]
    \item $\mathrm{e}^8$
    \item $\mathrm{e}^{15}$
    \item $2\mathrm{e}^8$
    \item $\mathrm{e}^{35}$
  \end{enumerate}

\item \textbf{[Q5]} Simplifier $\dfrac{\mathrm{e}^{7}}{\mathrm{e}^{3}}$.
  \begin{enumerate}[label=\Alph*.]
    \item $\mathrm{e}^{7/3}$
    \item $\mathrm{e}^{4}$
    \item $\mathrm{e}^{10}$
    \item $\mathrm{e}^{-21}$
  \end{enumerate}

\item \textbf{[Q6]} Simplifier $(\mathrm{e}^{2})^3 \times \mathrm{e}^{-4}$.
  \begin{enumerate}[label=\Alph*.]
    \item $\mathrm{e}^{-2}$
    \item $\mathrm{e}^{1}$
    \item $\mathrm{e}^{2}$
    \item $\mathrm{e}^{10}$
  \end{enumerate}

\item[] \textbf{Capacité C3}

\item \textbf{[Q7]} La fonction exponentielle est :
  \begin{enumerate}[label=\Alph*.]
    \item strictement décroissante sur $\mathbb{R}$
    \item constante sur $\mathbb{R}$
    \item croissante sur $]0\,;+\infty[$ et décroissante sur $]-\infty\,;0[$
    \item strictement croissante sur $\mathbb{R}$
  \end{enumerate}

\item \textbf{[Q8]} Pour un réel $x<0$, quelle affirmation situe correctement $\mathrm{e}^x$ ?
  \begin{enumerate}[label=\Alph*.]
    \item $0<\mathrm{e}^x<1$
    \item $\mathrm{e}^x=0$
    \item $\mathrm{e}^x<0$
    \item $\mathrm{e}^x>1$
  \end{enumerate}

\item \textbf{[Q9]} Sachant que la fonction exponentielle est strictement croissante, comparer $\mathrm{e}^{2{,}5}$ et $\mathrm{e}^{\sqrt{7}}$.
  \begin{enumerate}[label=\Alph*.]
    \item $\mathrm{e}^{2{,}5} > \mathrm{e}^{\sqrt{7}}$
    \item $\mathrm{e}^{2{,}5} < \mathrm{e}^{\sqrt{7}}$
    \item $\mathrm{e}^{2{,}5} = \mathrm{e}^{\sqrt{7}}$
    \item On ne peut pas comparer sans calculatrice.
  \end{enumerate}

\item[] \textbf{Capacité C4}

\item \textbf{[Q10]} Soit $f(t) = \mathrm{e}^{3t}$. Quelle est $f'(t)$ ?
  \begin{enumerate}[label=\Alph*.]
    \item $\mathrm{e}^{3t}$
    \item $3t\,\mathrm{e}^{3t}$
    \item $3\,\mathrm{e}^{3t}$
    \item $3\,\mathrm{e}^{3}$
  \end{enumerate}

\item \textbf{[Q11]} Soit $g(t) = \mathrm{e}^{-2t}$. Quelle est $g'(t)$ ?
  \begin{enumerate}[label=\Alph*.]
    \item $\mathrm{e}^{-2t}$
    \item $2\,\mathrm{e}^{-2t}$
    \item $-2\,\mathrm{e}^{2t}$
    \item $-2\,\mathrm{e}^{-2t}$
  \end{enumerate}

\item \textbf{[Q12]} Soit $h(t) = \mathrm{e}^{0{,}5t}$. Quelle est $h'(t)$ ?
  \begin{enumerate}[label=\Alph*.]
    \item $0{,}5\,\mathrm{e}^{0{,}5t}$
    \item $0{,}5t\,\mathrm{e}^{0{,}5t}$
    \item $2\,\mathrm{e}^{0{,}5t}$
    \item $\mathrm{e}^{0{,}5t}$
  \end{enumerate}

\item[] \textbf{Capacité C5}

\item \textbf{[Q13]} Une population est modélisée par $N(t)=500\,\mathrm{e}^{0{,}2t}$. Quelle est sa valeur initiale ?
  \begin{enumerate}[label=\Alph*.]
    \item $0$
    \item $500$
    \item $500\mathrm{e}^{0{,}2}$
    \item $0{,}2$
  \end{enumerate}

\item \textbf{[Q14]} Une masse est modélisée par $M(t)=80\,\mathrm{e}^{-0{,}3t}$. Quel est le quotient $M(t+1)/M(t)$ ?
  \begin{enumerate}[label=\Alph*.]
    \item $80\mathrm{e}^{-0{,}3}$
    \item $\mathrm{e}^{0{,}3}$
    \item $\mathrm{e}^{-0{,}3}$
    \item $-0{,}3$
  \end{enumerate}

\item \textbf{[Q15]} Pour $k>0$, quelle affirmation décrit la courbe de $t\mapsto\mathrm{e}^{-kt}$ ?
  \begin{enumerate}[label=\Alph*.]
    \item Elle est croissante et passe par $(0;0)$.
    \item Elle prend des valeurs négatives pour $t>0$.
    \item Elle est constante et égale a $1$.
    \item Elle est décroissante et passe par $(0;1)$.
  \end{enumerate}

\end{enumerate}

% NEXUS-QCM-TEACHER-ONLY-BEGIN
\ifnxVersionProfesseur
\clearpage
\section*{Clé de correction — réservée au professeur}

\begin{center}
\begin{tabular}{lll}
\hline
Question & Capacité & Réponse exacte \\
\hline
Q1 & C1 & \textbf{A} \\
Q2 & C1 & \textbf{B} \\
Q3 & C1 & \textbf{C} \\
Q4 & C2 & \textbf{A} \\
Q5 & C2 & \textbf{B} \\
Q6 & C2 & \textbf{C} \\
Q7 & C3 & \textbf{D} \\
Q8 & C3 & \textbf{A} \\
Q9 & C3 & \textbf{B} \\
Q10 & C4 & \textbf{C} \\
Q11 & C4 & \textbf{D} \\
Q12 & C4 & \textbf{A} \\
Q13 & C5 & \textbf{B} \\
Q14 & C5 & \textbf{C} \\
Q15 & C5 & \textbf{D} \\
\hline
\end{tabular}
\end{center}

\subsection*{Diagnostic des réponses erronées}

\begin{itemize}
  \item \textbf{Q1}
  \begin{itemize}
    \item \textbf{B} — La condition $f(1)=1$ définit la fonction $x\mapsto \mathrm{e}^{x-1}$, pas la fonction exponentielle $x\mapsto \mathrm{e}^x$, qui vérifie $f(0)=1$. \emph{Renvoi : M1, étape 1.}
    \item \textbf{C} — Avec $f(0) = 0$, la seule solution de $f' = f$ est la fonction nulle $f = 0$, pas l'exponentielle. La condition initiale $f(0) = 1$ est essentielle. \emph{Renvoi : M1, étape 1.}
    \item \textbf{D} — L'équation différentielle de l'exponentielle est $f' = f$, pas $f'(x) = xf(x)$. Confusion entre deux équations différentielles distinctes. \emph{Renvoi : M1, étape 1.}
  \end{itemize}
  \item \textbf{Q2}
  \begin{itemize}
    \item \textbf{A} — L'élève confond l'exposant $0$ avec le résultat : $\mathrm{e}^0 = 1$, pas $0$. Confusion entre puissance nulle et résultat nul. \emph{Renvoi : M1, propriété fondamentale.}
    \item \textbf{C} — L'élève confond $\mathrm{e}^0$ et $\mathrm{e}^1$. On a $\mathrm{e}^1 = \mathrm{e}$ mais $\mathrm{e}^0 = 1$. \emph{Renvoi : M1, propriété fondamentale.}
    \item \textbf{D} — L'exponentielle est définie sur $\mathbb{R}$ entier, donc $\mathrm{e}^0$ est bien défini et vaut $1$. \emph{Renvoi : M1, étape 1.}
  \end{itemize}
  \item \textbf{Q3}
  \begin{itemize}
    \item \textbf{A} — Pour $x < 0$, $\mathrm{e}^x$ est positif (par exemple $\mathrm{e}^{-1} = 1/\mathrm{e} \approx 0{,}37 > 0$). L'élève confond exposant négatif et résultat négatif. \emph{Renvoi : M1, étape 2.}
    \item \textbf{B} — L'exponentielle est strictement positive : $\mathrm{e}^x > 0$ pour tout $x$, elle n'atteint jamais $0$. \emph{Renvoi : M1, étape 2.}
    \item \textbf{D} — Pour $x < 0$, $\mathrm{e}^x < 1$ (par exemple $\mathrm{e}^{-1} \approx 0{,}37 < 1$). La propriété $\mathrm{e}^x > 1$ n'est vraie que pour $x > 0$. \emph{Renvoi : M1, étape 2 ; M3.}
  \end{itemize}
  \item \textbf{Q4}
  \begin{itemize}
    \item \textbf{B} — L'élève a multiplie les exposants ($3 \times 5 = 15$) au lieu de les additionner ($3 + 5 = 8$). Confusion entre $\mathrm{e}^a \times \mathrm{e}^b = \mathrm{e}^{a+b}$ et $(\mathrm{e}^a)^b = \mathrm{e}^{ab}$. \emph{Renvoi : M2, règle 1.}
    \item \textbf{C} — L'élève a bien additionne les exposants, mais a ensuite additionne a tort les coefficients implicites $1+1$. Dans un produit, ces coefficients se multiplient : $1\times1=1$. \emph{Renvoi : M2, règle 1 ; R4.}
    \item \textbf{D} — L'élève a concatene les chiffres $3$ et $5$ au lieu de les additionner. \emph{Renvoi : M2, règle 1 ; R4.}
  \end{itemize}
  \item \textbf{Q5}
  \begin{itemize}
    \item \textbf{A} — L'élève a divise les exposants ($7/3$) au lieu de les soustraire ($7 - 3 = 4$). Confusion entre quotient d'exponentielles et quotient d'exposants. \emph{Renvoi : M2, règle 2 ; R4.}
    \item \textbf{C} — L'élève a additionne les exposants ($7 + 3 = 10$) au lieu de les soustraire. Confusion entre produit et quotient d'exponentielles. \emph{Renvoi : M2, règle 2.}
    \item \textbf{D} — L'élève a multiplie les exposants avec un signe négatif ($-7 \times 3 = -21$) au lieu de soustraire. \emph{Renvoi : M2, règle 2 ; R4.}
  \end{itemize}
  \item \textbf{Q6}
  \begin{itemize}
    \item \textbf{A} — L'élève a oublie la puissance $3$ et a seulement calcule $\mathrm{e}^{2}\times\mathrm{e}^{-4}=\mathrm{e}^{-2}$. \emph{Renvoi : M2, règle 3 ; R4.}
    \item \textbf{B} — L'élève a calcule $(\mathrm{e}^2)^3 = \mathrm{e}^{2+3} = \mathrm{e}^{5}$ (addition au lieu de multiplication des exposants) puis $5 - 4 = 1$. \emph{Renvoi : M2, règle 3 ; R4.}
    \item \textbf{D} — L'élève a calcule $(\mathrm{e}^2)^3 = \mathrm{e}^6$ correctement, puis a additionne $6 + 4 = 10$ au lieu de soustraire $6 - 4 = 2$. \emph{Renvoi : M2, règles 2-3.}
  \end{itemize}
  \item \textbf{Q7}
  \begin{itemize}
    \item \textbf{A} — L'élève confond avec la fonction $x \mapsto \mathrm{e}^{-x}$ qui est décroissante. L'exponentielle $x \mapsto \mathrm{e}^x$ est strictement croissante car $\exp'(x) = \exp(x) > 0$. \emph{Renvoi : M3, étape 1.}
    \item \textbf{B} — La dérivée de l'exponentielle est $\exp(x) > 0$, donc la fonction est strictement croissante, certainement pas constante. \emph{Renvoi : M3, étape 1.}
    \item \textbf{C} — L'exponentielle n'a pas de minimum en $x = 0$. Elle est strictement croissante sur $\mathbb{R}$ tout entier, sans changement de variation. \emph{Renvoi : M3, étape 1.}
  \end{itemize}
  \item \textbf{Q8}
  \begin{itemize}
    \item \textbf{B} — La courbe de l'exponentielle ne coupe jamais l'axe des abscisses : $\mathrm{e}^x>0$. \emph{Renvoi : C3, courbe représentative.}
    \item \textbf{C} — Un exposant négatif ne rend pas la valeur négative : $\mathrm{e}^x$ est strictement positif pour tout réel $x$. \emph{Renvoi : C3, signe de l'exponentielle.}
    \item \textbf{D} — Comme $x<0$ et que l'exponentielle est croissante, $\mathrm{e}^x<\mathrm{e}^0=1$. \emph{Renvoi : C3, croissance de l'exponentielle.}
  \end{itemize}
  \item \textbf{Q9}
  \begin{itemize}
    \item \textbf{A} — L'élève a mal compare les exposants. Or $2{,}5^2 = 6{,}25 < 7$ donc $2{,}5 < \sqrt{7}$, et par croissance de l'exponentielle, $\mathrm{e}^{2{,}5} < \mathrm{e}^{\sqrt{7}}$. \emph{Renvoi : M3, étape 3.}
    \item \textbf{C} — $2{,}5 \neq \sqrt{7}$ car $2{,}5^2 = 6{,}25 \neq 7$. Comme les exposants sont différents et l'exponentielle est injective, les images sont différentes. \emph{Renvoi : M3 ; R3.}
    \item \textbf{D} — On peut comparer sans calculatrice en comparant les exposants : $2{,}5^2 = 6{,}25 < 7$ donc $2{,}5 < \sqrt{7}$, et on utilise la croissance de $\exp$. \emph{Renvoi : M3, étape 3.}
  \end{itemize}
  \item \textbf{Q10}
  \begin{itemize}
    \item \textbf{A} — L'élève a oublie le facteur constant $3$ : la dérivée de $t\mapsto\mathrm{e}^{at}$ est $t\mapsto a\mathrm{e}^{at}$. \emph{Renvoi : C4, dérivée de $t\mapsto\mathrm{e}^{at}$.}
    \item \textbf{B} — L'élève a multiplie par l'exposant $3t$ au lieu du coefficient constant $3$. \emph{Renvoi : C4, coefficient $a$.}
    \item \textbf{D} — Le facteur $3$ est correct, mais la fonction exponentielle conserve la variable : $\mathrm{e}^{3t}$ ne devient pas $\mathrm{e}^3$. \emph{Renvoi : C4, dérivée de $t\mapsto\mathrm{e}^{at}$.}
  \end{itemize}
  \item \textbf{Q11}
  \begin{itemize}
    \item \textbf{A} — L'élève a oublie le coefficient $-2$ qui multiplie $t$ dans l'exposant. \emph{Renvoi : C4, dérivée de $t\mapsto\mathrm{e}^{at}$.}
    \item \textbf{B} — Le coefficient a pour valeur $-2$, pas $2$ : le signe négatif a été perdu. \emph{Renvoi : C4, coefficient $a=-2$.}
    \item \textbf{C} — Le facteur $-2$ est correct, mais l'exposant reste $-2t$ : son signe ne change pas lors de la dérivation. \emph{Renvoi : C4, dérivée de $t\mapsto\mathrm{e}^{at}$.}
  \end{itemize}
  \item \textbf{Q12}
  \begin{itemize}
    \item \textbf{B} — L'élève a multiplie par l'exposant $0{,}5t$ au lieu du coefficient constant $0{,}5$. \emph{Renvoi : C4, coefficient $a$.}
    \item \textbf{C} — $2$ est l'inverse de $0{,}5$, pas le coefficient de $t$ dans l'exposant. \emph{Renvoi : C4, coefficient $a=0{,}5$.}
    \item \textbf{D} — L'élève a oublie le coefficient $0{,}5$ qui multiplie $t$ dans l'exposant. \emph{Renvoi : C4, coefficient $a=0{,}5$.}
  \end{itemize}
  \item \textbf{Q13}
  \begin{itemize}
    \item \textbf{A} — La valeur initiale est obtenue pour $t=0$ : le facteur exponentiel vaut alors $1$, pas $0$. \emph{Renvoi : C5, modèle exponentiel.}
    \item \textbf{C} — $500\mathrm{e}^{0{,}2}$ est la valeur au temps $t=1$, pas au temps initial $t=0$. \emph{Renvoi : C5, substitution de $t=0$.}
    \item \textbf{D} — $0{,}2$ est le taux continu du modèle ; ce n'est pas la valeur de la population. \emph{Renvoi : C5, paramètres du modèle.}
  \end{itemize}
  \item \textbf{Q14}
  \begin{itemize}
    \item \textbf{A} — Le coefficient initial $80$ se simplifie dans le quotient $M(t+1)/M(t)$. \emph{Renvoi : C5, coefficient multiplicateur.}
    \item \textbf{B} — Le signe de l'exposant a été inverse : le modèle de décroissance conserve le facteur $\mathrm{e}^{-0{,}3}$ d'une unité de temps a la suivante. \emph{Renvoi : C5, décroissance exponentielle.}
    \item \textbf{D} — $-0{,}3$ est le coefficient dans l'exposant ; le coefficient multiplicateur est son exponentielle $\mathrm{e}^{-0{,}3}$. \emph{Renvoi : C5, modèle exponentiel.}
  \end{itemize}
  \item \textbf{Q15}
  \begin{itemize}
    \item \textbf{A} — La courbe passe par $(0;1)$ car $\mathrm{e}^{0}=1$, et le coefficient $-k<0$ produit une décroissance. \emph{Renvoi : C5, courbe de $t\mapsto\mathrm{e}^{-kt}$.}
    \item \textbf{B} — Une exponentielle est toujours strictement positive, y compris lorsque son exposant est négatif. \emph{Renvoi : C3, signe ; C5, décroissance.}
    \item \textbf{C} — La fonction n'est constante que si le coefficient de $t$ dans l'exposant vaut $0$, or $-k<0$. \emph{Renvoi : C5, paramètre $k>0$.}
  \end{itemize}
\end{itemize}
\fi
% NEXUS-QCM-TEACHER-ONLY-END
